% SPDX-License-Identifier: CC-BY-4.0
% =====================================================================
% Section 1 of D. Joyner and Claude, "A hitch-hiker's guide to the
% metaplectic trace formula for SL(2)", revised for the 2-fold cover
% (n = 2), November 2026.
%
% New in the November 2026 draft, written in response to the second referee
% report (ChatGPT): overview, literature review, list of what is proved and
% what is cited, and the running example k = Q, S = {oo, 2}. It replaces the
% Introduction of ../original-1995/cctrace_new.tex for Sections 2-6; the
% stable form (the 1995 Section 7) is not yet revised.
% =====================================================================

\section{Introduction}\label{sec:intro}

Let $k$ be a number field, $G=\Res_{k/\Q}SL(2)$, and $\overline{G(\A)}$
the metaplectic double cover of $G(\A)=SL(2,\A_k)$. This paper writes out
the Arthur--Selberg trace formula for $\overline{G(\A)}$, first in its
non-invariant form and then in its invariant form, with every step
specialized to this group. The general theory exists: Arthur's trace
formula for connected reductive groups, and its extension to covering
groups by W.-W.~Li. The aim here is the opposite of generality. The group
has rank one, its only proper Levi subgroup is the diagonal torus $A$,
and the cover of $A$ is abelian. As a result most of the general
machinery reduces to statements that can be proved in a few lines and
checked numerically, and the terms of the formula can be written down.

\subsection{Overview}\label{sec:overview}

The paper has five parts, each resting on the ones before it.

\begin{description}
\item[The cover (Section~\ref{sec:background}).] The double cover is
  defined by Kubota's cocycle, place by place, and assembled into an
  adelic group that splits over $G(\Q)$. The cocycle is not continuous,
  so the topology of the cover and the spaces of test functions are
  defined intrinsically, and the coordinates $(g,\zeta)$ are used only as
  Borel coordinates. Test functions are \emph{anti-genuine}: they
  transform by $\zeta^{-1}$ under $\mu_2$.
\item[Rank $0$ (Section~\ref{sec:rank0}).] For the cover of a torus the
  trace formula is Poisson summation, twisted by a genuine automorphic
  character. For the split torus Weil's character $\chi_\psi$ is a
  canonical choice. This case is used twice later: to describe the
  cuspidal data of the Borel subgroup, and to match the correction terms
  on the two sides of the invariant formula.
\item[The geometric side (Section~\ref{sec:geometric}).] The kernel of
  $R(f)$ is truncated, and the integral of the truncated kernel is
  computed class by class: orbital integrals for the elliptic classes,
  weighted orbital integrals for the hyperbolic classes, and the finite
  part of a zeta integral for the two central classes.
\item[The spectral side (Section~\ref{sec:spectral}).] The same
  truncated integral is computed from the spectral decomposition. It
  gives the traces on the cuspidal spectrum and on the residual spectrum
  (even Weil representations), an integral of the logarithmic derivative
  of the intertwining operator $M(w,\lambda)$ along the unitary principal
  series, and one extra term $\frac14\operatorname{tr}(M(w,0)I(\mu_0,f))$
  for each character $\mu$ fixed by the Weyl group.
\item[The invariant form (Section~\ref{sec:invariant}).] The weighted
  orbital integrals and the term with $M^{-1}M'$ are not invariant under
  conjugation. Both fail in the same way, and subtracting a common
  correction, built from the normalized local intertwining operators by
  Fourier inversion on the cover of $A$, makes both sides invariant.
\end{description}

The result is Theorem~\ref{thm:itf} with Corollary~\ref{cor:itf-final}.
For $f=f_S\otimes e^S$, with $S$ a large enough finite set of places,
\begin{align*}
&\sum_{z=\pm1}\Bigl(\mathrm{vol}\cdot f_S(z)+\sum_{\xi}a^G(S,u_\xi)\,
J_G(zu_\xi,f_S)+\tfrac12c_k\,I_A(z,f_S)\Bigr)\\
&\qquad+\sum_{\mathcal O\ \mathrm{elliptic}}\mathrm{vol}(T_\mu)\,J_G(\mu_S,f_S)
+\sum_{\mathcal O\ \mathrm{hyperbolic}}c_k\,I_A(\gamma_S,f_S)\\
&=\sum_\pi m_{\mathrm{disc}}(\pi)\operatorname{tr}\pi(f)
-\frac14\sum_{b\in k^\times/k^{\times2}}\operatorname{tr}I(\chi_{\psi_b,0},f)\\
&\qquad-\frac12\sum_\mu\frac1{2\pi}\int_\R\frac{m'}m(\mu,it)\,\operatorname{tr}
I(\mu_{it},f)\,dt .
\end{align*}
Here $J_G$ are orbital integrals, $I_A$ the invariant distributions that
replace the weighted orbital integrals, $\pi$ runs over the discrete
spectrum (cuspidal and residual), $\mu$ over the genuine automorphic
characters of the cover of $A$, and $m$ is the global normalizing factor
of the intertwining operator, essentially
$L(2\lambda,\chi^2)/L(2\lambda+1,\chi^2)$. The notation is collected in
the Glossary at the end.

Three features distinguish the double cover from $SL(2)$ itself, and they
are the main mathematical content of Sections~\ref{sec:spectral}
and~\ref{sec:invariant}. First, the Weyl group does not act on the cover
of $A$ by inversion; the characters it fixes are the Weil characters
$\chi_{\psi_b}$, one for each square class $b$
(Lemma~\ref{lem:w-action}). Second, the residual spectrum is not the
trivial representation but the sum of the even Weil representations
$r^+_{\psi_{-b}}$, which come from the pole of $M(w,\lambda)$ at
$\lambda=\frac12$ (Theorem~\ref{thm:residual}). Third, $M(w,0)=-1$ for
every one of these characters (Proposition~\ref{prop:Mw0}), whereas for
$SL(2)$ and a nontrivial quadratic character the corresponding operator
is not a scalar and is the source of the multiplicity formula of Labesse
and Langlands~\cite{LL}.

\subsection{Literature}\label{sec:literature}

\paragraph{Covering groups and the Weil representation.} The double
cover of $SL(2)$ over a local field and over the adeles first appeared in
Weil's work on the representation that now carries his name~\cite{We64}.
Kubota gave the explicit cocycle used here, for the $n$-fold cover of
$SL(2)$ over a local field containing the $n$-th roots of unity, and the
splitting over the rational points~\cite{K1,K2}. Moore~\cite{M} and
Matsumoto~\cite{Mat} treated the central extensions of split groups
over local fields and adele rings in general. Kazhdan and
Patterson~\cite{KP} studied the covers of $GL(r)$ and their automorphic
forms, and Brylinski and Deligne~\cite{BD} placed all these covers in one
algebraic framework, as extensions by $K_2$. Gelbart's lecture
notes~\cite{Ge2} remain the standard account of the double cover of
$SL(2)$ and $GL(2)$, the Weil representation and theta functions.
Covering groups are not algebraic groups, and the Borel cocycles that
describe them are not continuous; the relevant general theorem on the
topology of such extensions is Mackey's~\cite{Ma57}. For the Langlands
program for covering groups, and its history, see Gan, Gao and
Weissman~\cite{GGW} and Weissman~\cite{We18}.

\paragraph{Half-integral weight and the Shimura correspondence.}
Shimura~\cite{Sh} attached to a cusp form of weight $k/2$ a form of even
weight $k-1$. Niwa~\cite{Ni} and Shintani~\cite{Shi} realized the
correspondence and its inverse by theta series. Gelbart and
Piatetski-Shapiro~\cite{GPS80,GPS83} and Waldspurger~\cite{Wa80,Wa81,Wa91}
recast it in terms of automorphic representations of the double cover.
Waldspurger's description of the discrete spectrum of the double cover of
$SL(2)$, through the theta correspondence with $PGL(2)$ and its inner
forms, is the definitive result for the group studied here. The
identification of the residual spectrum with Weil representations, used
in Theorem~\ref{thm:residual}, belongs to this circle of
results~\cite{Ge2,GPS80,Wa80}. Gan and Savin~\cite{GS12}
obtained the local Langlands correspondence for the double cover of
$Sp(2n)$ from the theta correspondence, and Gan and Ichino~\cite{GI}
extended Waldspurger's global description to the generic part of the
discrete spectrum of the double cover of $Sp(2n)$.

\paragraph{The trace formula.} The trace formula goes back to
Selberg~\cite{Sel}. For $GL(2)$ it was put in adelic form by Jacquet and
Langlands~\cite{JL}; the treatment of Gelbart and Jacquet~\cite{GJ} is the
closest in spirit to the present paper, and we compare signs and
constants with it in Section~\ref{sec:spectral}. The classical Selberg
trace formula with multiplier systems and real weight, as in
Hejhal~\cite{He}, covers forms of half-integral weight on congruence
subgroups of $SL(2,\Z)$. For a general connected reductive group the
formula is Arthur's: the coarse expansions~\cite{A2,A3}, the invariant
form~\cite{A6}, the fine expansions~\cite{A12,A14}, the local theory of
weighted orbital integrals and normalized intertwining
operators~\cite{A7,A9,A11}, and the invariant trace
formula~\cite{A10}; his Clay lectures~\cite{Ar05} are the best
introduction. The spectral decomposition on which the spectral side
rests is due to Langlands~\cite{L2}, and M\oe glin and
Waldspurger~\cite{MW2} wrote it in a form that includes finite central
covers. For $SL(2)$ the stable trace formula is the subject of Labesse
and Langlands~\cite{LL}.

\paragraph{The trace formula for covering groups.} Flicker~\cite{Fl} used
a trace formula for the $n$-fold covers of $GL(2)$ to prove the
metaplectic correspondence with $GL(2)$, and Flicker and
Kazhdan~\cite{FK} did the same for $GL(r)$; see also Kazhdan and Patterson~\cite{KP2}. Mezo~\cite{Me1,Me2} prepared
the comparison of the full invariant trace formulas of $GL(r)$ and its
covers. The first author studied a metaplectic analog of Kazhdan's
endoscopic lifting in~\cite{J88}. The 1995 version of the present notes
wrote out Arthur's formula for the $n$-fold cover of $SL(2)$, with
several steps conjectural. Those steps are now theorems of
Li, who extended Arthur's trace formula to covering groups of connected
reductive groups in the four papers~\cite{LiI,LiII,LiIII,LiIV}: the fine
geometric expansion, the local harmonic analysis including normalized
intertwining operators, the fine spectral expansion, and the invariant
distributions (the invariant formula under a hypothesis at the real
places, which he verifies for the double cover of $Sp(2n)$). For the double cover of $Sp(2n)$ Li went on to
develop endoscopy, building on Adams~\cite{Ad} and Renard~\cite{Re} at
the real place: the transfer of orbital integrals~\cite{Li11}, the
weighted fundamental lemma~\cite{Li12}, the stabilization of the elliptic
terms~\cite{Li15}, local character relations~\cite{Li19}, the
stabilization of the full trace formula~\cite{Li21}, and a multiplicity
formula for the discrete spectrum in terms of Arthur
parameters~\cite{Li24}. For $n=1$ these results apply to the group of
this paper (Li works with an eight-fold cover that contains the double
cover). A $p$-adic character lifting between
the double cover of $SL(2)$ and $SO(1,2)$ is in Schultz's
thesis~\cite{Sch}.

\paragraph{Local harmonic analysis on the double cover of $SL(2)$.} The
irreducibility of the unitary genuine principal series is due to
Moen~\cite{Mo2}. Szpruch~\cite{Sz09} computed the local coefficients of
the genuine principal series of the double cover of $SL(2,F)$, and
Goldberg and Szpruch~\cite{GS16} the Plancherel measures for the $n$-fold
covers. These results bear on the normalizing factors at the places where
our Section~\ref{sec:normalized} is not explicit. The unramified computation
of the intertwining operator (the Gindikin--Karpelevich formula) for
general covers is in~\cite{KP,McN,Gao}; see also Savin~\cite{Sav} for the
Hecke algebra correspondence.

\paragraph{The place of this paper.} Nothing in
Sections~\ref{sec:background}--\ref{sec:invariant} goes beyond what the
general theorems of Arthur and Li give. What the paper adds is an
account in which the group is fixed, the constants are determined, the
proofs that are short for this group are given in full, and the
remaining appeals to the literature are listed. We hope that it is
useful as a worked example of Li's theory and as a companion to the
classical treatments of half-integral weight.

\subsection{What is proved here and what is cited}\label{sec:proved-cited}

Each section ends with a subsection ``Status of the results'', which
lists the statements taken from the literature and the statements
checked numerically. In summary:

\begin{center}
\begin{tabular}{p{0.11\textwidth}p{0.40\textwidth}p{0.40\textwidth}}
Section & Proved in the text & Taken from the literature\\\hline
\ref{sec:background} & the splittings, the center, commutativity of
  preimages of tori, the local triviality of the cocycle, the adelic
  cover and its splitting over $G(\Q)$ &
  Kubota's cocycle and the existence of the topological
  cover~\cite{K1,M}; the splitting over $K_v$~\cite{K2}; strong
  approximation\\
\ref{sec:rank0} & the trace formula for covers of tori and its
  reduction to Poisson summation & the properties of the Weil
  index~\cite{We64}\\
\ref{sec:geometric} & the coarse geometric expansion and the
  polynomials $J^T_{\mathcal O}$; the local distribution $J_A(z,f)$ &
  the fine expansion at the central classes~\cite{A12,LiI}, whose
  coefficients $a^G(S,u_\xi)$ are not evaluated\\
\ref{sec:spectral} & cuspidal data and the Weyl action; the unramified
  intertwining operator; the terms $J^T_\chi$ from the Maass--Selberg
  relations; the matching $V_b\cong r^+_{\psi_{-b}}$ &
  the spectral decomposition and Eisenstein series~\cite{L2,MW2}; the
  convergence and polynomiality of the spectral
  side~\cite{A3,LiIII}; the residual spectrum as theta
  functions~\cite{Ge2,GPS80,Wa80}\\
\ref{sec:invariant} & the variance formulas; the real normalizing
  factors; the invariance of $I_A$; the invariant trace formula;
  $M(w,0)=-1$ & the existence of normalizing factors~\cite{A11,LiII},
  which are not explicit at the dyadic and ramified places
\end{tabular}
\end{center}

Two ingredients of the final formula are therefore defined but not
evaluated: the unipotent coefficients $a^G(S,u_\xi)$, and the local
normalizing factors $r_v$ at the places where the additive character or
the inducing character is ramified, and at the dyadic places. Remarks
numbered ``2026.$n$'' record these and other points that rest on
citation or on numerical checks only. The explicit formulas are tested
by the scripts in \texttt{checks/}; they are evidence, not proofs.

The stable form of the trace formula, Section~7 of the 1995 notes, is not
part of this draft. The 1995 notes define stable conjugacy on the cover
through conjugation in a cover of $GL(2)$, following~\cite{LL}; its
relation to Li's endoscopy~\cite{Li11,Li15} has to be worked out when
that section is revised. Section~\ref{sec:gl2} prepares for it.

\subsection{A running example}\label{sec:running}

The case $k=\Q$, $S=\{\infty,2\}$ is the smallest one in which every
feature of the formula appears, and it is followed through the paper. We
collect it here; $\psi$ is the character of $\A/\Q$ with
$\psi_\infty(x)=e^{-2\pi ix}$.

\begin{itemize}
\item \emph{The cover.} At the real place the cover is $Mp(2,\R)$; its
  center is cyclic of order $4$, and so is the center at the place $2$
  (Lemma~\ref{lem:local}(e), Example~\ref{ex:real}). At an odd prime $p$
  the cover splits over $SL(2,\Z_p)$ by $k\mapsto(k,\kappa_p(k))$, and on
  $SL(2,\Z)$ the product of the $\kappa_p$ is a Jacobi symbol
  (Example~\ref{ex:jacobi}). At $2$ we use a splitting over a congruence
  subgroup (Lemma~\ref{lem:local}(c), Remark~\ref{rem:dyadic}). Test functions are
  $f=f_\infty\otimes f_2\otimes\bigotimes_{p\ \mathrm{odd}}e_p$.
\item \emph{Rank $0$.} $\chi_{\psi,\infty}(a(-1),1)=i$ and
  $\chi_{\psi,p}(a(p),1)\in\{1,i\}$ for odd $p$
  (Examples~\ref{ex:order4} and~\ref{ex:weil-Q}): a genuine character of
  the cover of $A$ is not an ordinary character times a sign.
\item \emph{Geometric side.} $c_\Q=1$. Besides the elliptic classes
  there are the two central classes $\pm1$ and the hyperbolic classes of
  $a(t)$, $t=\pm2^j$, $j\ge1$. The unipotent terms are indexed by
  $\Q_S^\times/\Q_S^{\times2}$, which has $2\cdot8=16$ elements. The
  constant term of the unipotent contribution is computed for a Gaussian
  in Example~\ref{ex:unip-Q}.
\item \emph{Spectral side.} The characters fixed by the Weyl group that
  are unramified outside $S$ are the four $\chi_{\psi_b}$,
  $b\in\{\pm1,\pm2\}$ (Example~\ref{ex:ramified-Q}). Each carries a
  residual representation $V_b\cong r^+_{\psi_{-b}}$; $V_1$ contains the
  classical theta function $\sum_ne^{2\pi in^2z}$
  (Example~\ref{ex:theta}). The partial normalizing factor is
  $(1-2^{-2\lambda})\zeta(2\lambda)/((1-2^{-1-2\lambda})\zeta(2\lambda+1))$,
  with residue $2/\pi^2$ at $\lambda=\frac12$
  (Examples~\ref{ex:theta} and~\ref{ex:mprime-Q}).
\item \emph{Invariant form.} At the real place the normalized
  intertwining operator is $1$ on the $K$-type of weight $\pm\frac12$ and
  a product of factors $(c-\lambda)/(c+\lambda)$ on the others
  (Example~\ref{ex:real-R}), and the correction $\Phi_A$ is an Abel
  transform smoothed by an exponentially decaying kernel
  (Example~\ref{ex:Phi-real}). The only bad place is $2$, where the
  normalizing factor is not explicit. The discrete correction is
  $-\frac14\sum_{b=\pm1,\pm2}\operatorname{tr}I(\chi_{\psi_b,0},f)$
  (Example~\ref{ex:itf-Q}).
\end{itemize}

\subsection{Conventions}\label{sec:conventions-intro}

The revision treats the double cover only; the 1995 notes treated the
$n$-fold cover, and remarks in Sections~\ref{sec:background}
and~\ref{sec:rank0} say what changes for $n>2$. Comparisons with the
1995 notes and with the 1997 book manuscript, including the corrections
made, are collected in Appendix~\ref{sec:appendix-1995} and are not
needed to read the paper. Notation and terms are listed in the Glossary.
